%
% claude-code-quarry.tex
%
% Z Specification: Claude Code + Quarry Knowledge Capture State (Layer 2)
% Extends the base Claude Code state machine with quarry automatic
% knowledge capture, codebase indexing, and compaction safety.
%
% Type-check: fuzz -t claude-code-quarry.tex
% Animate:    probcli claude-code-quarry.tex -animate 20
%

\documentclass[a4paper,10pt,fleqn]{article}
\usepackage[margin=1in]{geometry}
\usepackage{fuzz}
\usepackage[colorlinks=true,linkcolor=blue,citecolor=blue,urlcolor=blue]{hyperref}

\begin{document}

\title{Claude Code + Quarry: A Z Specification}
\author{Layer 2 --- Knowledge Capture State}
\date{March 2026}
\hypersetup{
  pdftitle={Claude Code + Quarry: A Z Specification},
  pdfauthor={Punt Labs},
  pdfsubject={Z Specification},
  pdfcreator={fuzz/probcli}
}
\maketitle

\tableofcontents
\newpage

% ============================================================
\section{Introduction}
% ============================================================

This specification extends the base Claude Code state machine
(\texttt{claude-code.tex}) with quarry automatic knowledge capture
state. Quarry provides local semantic search over codebases,
documents, and conversation history.

The key properties verified by this model:

\begin{enumerate}
  \item \textbf{Knowledge before work}: the codebase is indexed
    before Claude begins processing. The \texttt{SessionStart} hook
    auto-registers and syncs the working directory.
  \item \textbf{No knowledge loss on compaction}: before context
    compaction shrinks the conversation, the \texttt{PreCompact} hook
    captures session notes. No path reaches \texttt{CompactContext}
    without the capture flag being set.
  \item \textbf{WebFetch auto-capture}: URLs fetched during a session
    are automatically ingested into the \texttt{web-captures} collection.
  \item \textbf{Capture toggles are respected}: each auto-capture
    layer (session sync, web fetch, compaction) can be independently
    disabled via \texttt{.claude/quarry.local.md}.
\end{enumerate}

% ============================================================
\section{Base Types (from claude-code.tex)}
% ============================================================

\begin{zed}
  [TOOLNAME, TOOLINPUT, TOOLOUTPUT, PROMPT, RESPONSE, AGENTID, SESSIONID, CONTEXTCHUNK]
\end{zed}

\begin{zed}
  SessionPhase ::= spInactive | spIdle | spProcessing | spResponding | spEnded
\end{zed}

\begin{zed}
  ToolPhase ::= tpNone | tpPreHook | tpPermission | tpExecuting | tpPostHook
\end{zed}

\begin{zed}
  PermissionMode ::= pmDefault | pmPlan | pmAcceptEdits | pmDontAsk | pmBypass
\end{zed}

\begin{zed}
  PermissionDecision ::= pdAllow | pdDeny | pdAsk
\end{zed}

\begin{zed}
  HookDecision ::= hdAllow | hdBlock
\end{zed}

\begin{zed}
  SessionStartSource ::= ssStartup | ssResume | ssClear | ssCompact
\end{zed}

\begin{zed}
  SessionEndReason ::= seClear | seLogout | sePromptExit | seBypassDisabled | seOther
\end{zed}

\begin{zed}
  ZBOOL ::= ztrue | zfalse
\end{zed}

% ============================================================
\section{Quarry Types}
% ============================================================

\begin{zed}
  [DBNAME, URL]
\end{zed}

\texttt{DBNAME} identifies a quarry database (e.g., \texttt{default},
\texttt{punt-labs}). \texttt{URL} identifies a fetched web resource.

% ============================================================
\section{Global Constants}
% ============================================================

\begin{axdef}
  maxTools : \nat \\
  maxSubagents : \nat \\
  maxContext : \nat \\
  maxCaptures : \nat
  \where
  maxTools = 5 \\
  maxSubagents = 2 \\
  maxContext = 3 \\
  maxCaptures = 5
\end{axdef}

% ============================================================
\section{State}
% ============================================================

\begin{schema}{State}
  %
  % --- Base Claude Code state ---
  %
  sessionId : \finset SESSIONID \\
  sessionPhase : SessionPhase \\
  permMode : PermissionMode \\
  toolPhase : ToolPhase \\
  currentTool : \finset TOOLNAME \\
  toolsThisTurn : \nat \\
  activeSubagents : \finset AGENTID \\
  context : \seq CONTEXTCHUNK \\
  pendingPrompt : \finset PROMPT \\
  pendingResponse : \finset RESPONSE \\
  %
  % --- Quarry knowledge capture state ---
  %
  codebaseIndexed : ZBOOL \\
  activeDb : \finset DBNAME \\
  capturedUrls : \finset URL \\
  compactionSaved : ZBOOL \\
  %
  % --- Capture config (from .claude/quarry.local.md) ---
  %
  cfgSessionSync : ZBOOL \\
  cfgWebFetch : ZBOOL \\
  cfgCompaction : ZBOOL
  \where
  %
  % === Base invariants ===
  %
  \# sessionId \leq 1 \\
  sessionPhase = spInactive \implies sessionId = \emptyset \\
  sessionPhase \neq spInactive \implies \# sessionId = 1 \\
  sessionPhase \neq spProcessing \implies toolPhase = tpNone \\
  \# currentTool \leq 1 \\
  toolPhase = tpNone \implies currentTool = \emptyset \\
  toolPhase \neq tpNone \implies \# currentTool = 1 \\
  toolsThisTurn \leq maxTools \\
  sessionPhase = spInactive \implies activeSubagents = \emptyset \\
  sessionPhase = spEnded \implies activeSubagents = \emptyset \\
  \# activeSubagents \leq maxSubagents \\
  \# context \leq maxContext \\
  \# pendingPrompt \leq 1 \\
  sessionPhase = spIdle \implies pendingPrompt = \emptyset \\
  sessionPhase = spInactive \implies pendingPrompt = \emptyset \\
  sessionPhase = spEnded \implies pendingPrompt = \emptyset \\
  \# pendingResponse \leq 1 \\
  sessionPhase \neq spResponding \implies pendingResponse = \emptyset \\
  %
  % === Quarry invariants ===
  %
  % Active database: at most one
  %
  \# activeDb \leq 1 \\
  %
  % Captured URLs bounded for animation
  %
  \# capturedUrls \leq maxCaptures \\
  %
  % No quarry state when session is inactive
  %
  sessionPhase = spInactive \implies codebaseIndexed = zfalse \\
  sessionPhase = spInactive \implies activeDb = \emptyset \\
  sessionPhase = spInactive \implies capturedUrls = \emptyset \\
  sessionPhase = spInactive \implies compactionSaved = zfalse
\end{schema}

% ============================================================
\section{Initialization}
% ============================================================

\texttt{Init} describes the machine before any session exists. The three
capture toggles are read from \texttt{.claude/quarry.local.md} by the
\texttt{SessionStart} hook, so at this point they have not been read and
their values mean nothing; they are pinned to \texttt{zfalse} for
definiteness, alongside the other quarry fields, all of which are empty
or false while \texttt{sessionPhase = spInactive}.

Setting the configuration here instead would assert that the machine
knows it before the hook that reads it has run. That is the error this
document already made in a subtler form: the flags were fixed to
\texttt{ztrue} at initialisation and carried forward unchanged by every
operation, so the configuration was declared as state and then made
constant. The fourth key property promised above --- that each capture
layer can be independently disabled --- was stated over flags no
execution could vary, and the operations guarded by \texttt{zfalse} were
unreachable from the initial state. A specification that documents
behaviour it cannot exhibit is worse than one that omits it, because it
reads as though it had been checked.

The configuration therefore enters where it is read, as inputs on the two
\texttt{SessionStart} operations below. It remains immutable for the life
of the session: \texttt{QuarryFrame} carries \texttt{cfg\ldots' =
cfg\ldots} through every operation that follows.

\begin{schema}{Init}
  State'
  \where
  sessionId' = \emptyset \\
  sessionPhase' = spInactive \\
  permMode' = pmDefault \\
  toolPhase' = tpNone \\
  currentTool' = \emptyset \\
  toolsThisTurn' = 0 \\
  activeSubagents' = \emptyset \\
  context' = \langle \rangle \\
  pendingPrompt' = \emptyset \\
  pendingResponse' = \emptyset \\
  %
  codebaseIndexed' = zfalse \\
  activeDb' = \emptyset \\
  capturedUrls' = \emptyset \\
  compactionSaved' = zfalse \\
  cfgSessionSync' = zfalse \\
  cfgWebFetch' = zfalse \\
  cfgCompaction' = zfalse
\end{schema}

% ============================================================
\section{Quarry Frame Condition}
% ============================================================

\begin{schema}{QuarryFrame}
  \Delta State
  \where
  codebaseIndexed' = codebaseIndexed \\
  activeDb' = activeDb \\
  capturedUrls' = capturedUrls \\
  compactionSaved' = compactionSaved \\
  cfgSessionSync' = cfgSessionSync \\
  cfgWebFetch' = cfgWebFetch \\
  cfgCompaction' = cfgCompaction
\end{schema}

% ============================================================
\section{Session Lifecycle Operations}
% ============================================================

\subsection{StartSessionWithSync}

The \texttt{SessionStart} hook fires and reads the capture configuration.
This is the branch it takes when session sync is enabled: quarry
auto-registers the working directory and syncs, so the codebase is
indexed before Claude can process a prompt.

The two \texttt{SessionStart} schemas have the same precondition and
differ in what they establish. That is deliberate: which one fires is
decided by the configuration file, and a nondeterministic choice between
two operations is how this specification says that something outside the
machine decides. Guarding them on \texttt{cfgSessionSync} instead would
require the flag to be known before the hook that reads it has run.

\begin{schema}{StartSessionWithSync}
  \Delta State \\
  sid? : SESSIONID \\
  source? : SessionStartSource \\
  mode? : PermissionMode \\
  webFetch? : ZBOOL \\
  compaction? : ZBOOL
  \where
  sessionPhase = spInactive \\
  %
  sessionPhase' = spIdle \\
  sessionId' = \{ sid? \} \\
  permMode' = mode? \\
  toolPhase' = tpNone \\
  currentTool' = \emptyset \\
  toolsThisTurn' = 0 \\
  activeSubagents' = \emptyset \\
  context' = \langle \rangle \\
  pendingPrompt' = \emptyset \\
  pendingResponse' = \emptyset \\
  %
  codebaseIndexed' = ztrue \\
  activeDb' = \emptyset \\
  capturedUrls' = \emptyset \\
  compactionSaved' = zfalse \\
  cfgSessionSync' = ztrue \\
  cfgWebFetch' = webFetch? \\
  cfgCompaction' = compaction?
\end{schema}

\subsection{StartSessionNoSync}

The other branch: the session starts with sync disabled, and the codebase
is not indexed automatically.

\begin{schema}{StartSessionNoSync}
  \Delta State \\
  sid? : SESSIONID \\
  source? : SessionStartSource \\
  mode? : PermissionMode \\
  webFetch? : ZBOOL \\
  compaction? : ZBOOL
  \where
  sessionPhase = spInactive \\
  %
  sessionPhase' = spIdle \\
  sessionId' = \{ sid? \} \\
  permMode' = mode? \\
  toolPhase' = tpNone \\
  currentTool' = \emptyset \\
  toolsThisTurn' = 0 \\
  activeSubagents' = \emptyset \\
  context' = \langle \rangle \\
  pendingPrompt' = \emptyset \\
  pendingResponse' = \emptyset \\
  %
  codebaseIndexed' = zfalse \\
  activeDb' = \emptyset \\
  capturedUrls' = \emptyset \\
  compactionSaved' = zfalse \\
  cfgSessionSync' = zfalse \\
  cfgWebFetch' = webFetch? \\
  cfgCompaction' = compaction?
\end{schema}

\subsection{SubmitPrompt}

\begin{schema}{SubmitPrompt}
  QuarryFrame \\
  prompt? : PROMPT \\
  hookResult? : HookDecision \\
  contextChunk? : CONTEXTCHUNK
  \where
  sessionPhase = spIdle \\
  hookResult? = hdAllow \\
  sessionPhase' = spProcessing \\
  pendingPrompt' = \{ prompt? \} \\
  context' = (1 \upto maxContext) \dres (context \cat \langle contextChunk? \rangle) \\
  toolsThisTurn' = 0 \\
  toolPhase' = tpNone \\
  currentTool' = \emptyset \\
  sessionId' = sessionId \\
  permMode' = permMode \\
  activeSubagents' = activeSubagents \\
  pendingResponse' = \emptyset
\end{schema}

\begin{schema}{RejectPrompt}
  \Xi State \\
  prompt? : PROMPT \\
  hookResult? : HookDecision
  \where
  sessionPhase = spIdle \\
  hookResult? = hdBlock
\end{schema}

% ============================================================
\section{Tool Call Operations}
% ============================================================

\subsection{BeginToolCall}

\begin{schema}{BeginToolCall}
  QuarryFrame \\
  tool? : TOOLNAME
  \where
  sessionPhase = spProcessing \\
  toolPhase = tpNone \\
  toolsThisTurn < maxTools \\
  toolPhase' = tpPreHook \\
  currentTool' = \{ tool? \} \\
  sessionPhase' = sessionPhase \\
  sessionId' = sessionId \\
  permMode' = permMode \\
  toolsThisTurn' = toolsThisTurn \\
  activeSubagents' = activeSubagents \\
  context' = context \\
  pendingPrompt' = pendingPrompt \\
  pendingResponse' = pendingResponse
\end{schema}

\subsection{PreToolHookAllow}

\begin{schema}{PreToolHookAllow}
  QuarryFrame \\
  hookDecision? : PermissionDecision \\
  needsPermission? : ZBOOL
  \where
  sessionPhase = spProcessing \\
  toolPhase = tpPreHook \\
  hookDecision? = pdAllow \\
  toolPhase' = \IF needsPermission? = ztrue \THEN tpPermission \ELSE tpExecuting \\
  sessionPhase' = sessionPhase \\
  sessionId' = sessionId \\
  permMode' = permMode \\
  currentTool' = currentTool \\
  toolsThisTurn' = toolsThisTurn \\
  activeSubagents' = activeSubagents \\
  context' = context \\
  pendingPrompt' = pendingPrompt \\
  pendingResponse' = pendingResponse
\end{schema}

\subsection{PreToolHookDeny}

\begin{schema}{PreToolHookDeny}
  QuarryFrame \\
  hookDecision? : PermissionDecision \\
  denialContext? : CONTEXTCHUNK
  \where
  sessionPhase = spProcessing \\
  toolPhase = tpPreHook \\
  hookDecision? = pdDeny \\
  toolPhase' = tpNone \\
  currentTool' = \emptyset \\
  context' = (1 \upto maxContext) \dres (context \cat \langle denialContext? \rangle) \\
  sessionPhase' = sessionPhase \\
  sessionId' = sessionId \\
  permMode' = permMode \\
  toolsThisTurn' = toolsThisTurn \\
  activeSubagents' = activeSubagents \\
  pendingPrompt' = pendingPrompt \\
  pendingResponse' = pendingResponse
\end{schema}

\subsection{GrantPermission}

\begin{schema}{GrantPermission}
  QuarryFrame
  \where
  sessionPhase = spProcessing \\
  toolPhase = tpPermission \\
  toolPhase' = tpExecuting \\
  sessionPhase' = sessionPhase \\
  sessionId' = sessionId \\
  permMode' = permMode \\
  currentTool' = currentTool \\
  toolsThisTurn' = toolsThisTurn \\
  activeSubagents' = activeSubagents \\
  context' = context \\
  pendingPrompt' = pendingPrompt \\
  pendingResponse' = pendingResponse
\end{schema}

\subsection{DenyPermission}

\begin{schema}{DenyPermission}
  QuarryFrame \\
  denialContext? : CONTEXTCHUNK
  \where
  sessionPhase = spProcessing \\
  toolPhase = tpPermission \\
  toolPhase' = tpNone \\
  currentTool' = \emptyset \\
  context' = (1 \upto maxContext) \dres (context \cat \langle denialContext? \rangle) \\
  sessionPhase' = sessionPhase \\
  sessionId' = sessionId \\
  permMode' = permMode \\
  toolsThisTurn' = toolsThisTurn \\
  activeSubagents' = activeSubagents \\
  pendingPrompt' = pendingPrompt \\
  pendingResponse' = pendingResponse
\end{schema}

\subsection{CompleteToolSuccess}

\begin{schema}{CompleteToolSuccess}
  QuarryFrame \\
  output? : CONTEXTCHUNK
  \where
  sessionPhase = spProcessing \\
  toolPhase = tpExecuting \\
  toolPhase' = tpNone \\
  currentTool' = \emptyset \\
  toolsThisTurn' = toolsThisTurn + 1 \\
  context' = (1 \upto maxContext) \dres (context \cat \langle output? \rangle) \\
  sessionPhase' = sessionPhase \\
  sessionId' = sessionId \\
  permMode' = permMode \\
  activeSubagents' = activeSubagents \\
  pendingPrompt' = pendingPrompt \\
  pendingResponse' = pendingResponse
\end{schema}

\subsection{CompleteToolFailure}

\begin{schema}{CompleteToolFailure}
  QuarryFrame \\
  errorContext? : CONTEXTCHUNK
  \where
  sessionPhase = spProcessing \\
  toolPhase = tpExecuting \\
  toolPhase' = tpNone \\
  currentTool' = \emptyset \\
  toolsThisTurn' = toolsThisTurn + 1 \\
  context' = (1 \upto maxContext) \dres (context \cat \langle errorContext? \rangle) \\
  sessionPhase' = sessionPhase \\
  sessionId' = sessionId \\
  permMode' = permMode \\
  activeSubagents' = activeSubagents \\
  pendingPrompt' = pendingPrompt \\
  pendingResponse' = pendingResponse
\end{schema}

% ============================================================
\section{Quarry Operations}
% ============================================================

\subsection{CaptureWebFetch}

\texttt{PostToolUse} on \texttt{WebFetch}. When web fetch capture
is enabled, the fetched URL is added to the captured set and
ingested into the \texttt{web-captures} collection (async,
fire-and-forget).

\begin{schema}{CaptureWebFetch}
  \Delta State \\
  url? : URL
  \where
  sessionPhase = spProcessing \\
  toolPhase = tpNone \\
  cfgWebFetch = ztrue \\
  url? \notin capturedUrls \\
  \# capturedUrls < maxCaptures \\
  %
  capturedUrls' = capturedUrls \cup \{ url? \} \\
  %
  sessionPhase' = sessionPhase \\
  sessionId' = sessionId \\
  permMode' = permMode \\
  toolPhase' = toolPhase \\
  currentTool' = currentTool \\
  toolsThisTurn' = toolsThisTurn \\
  activeSubagents' = activeSubagents \\
  context' = context \\
  pendingPrompt' = pendingPrompt \\
  pendingResponse' = pendingResponse \\
  codebaseIndexed' = codebaseIndexed \\
  activeDb' = activeDb \\
  compactionSaved' = compactionSaved \\
  cfgSessionSync' = cfgSessionSync \\
  cfgWebFetch' = cfgWebFetch \\
  cfgCompaction' = cfgCompaction
\end{schema}

\subsection{SwitchDatabase}

Claude calls \texttt{quarry use <name>}. Changes the active database
for subsequent queries.

\begin{schema}{SwitchDatabase}
  \Delta State \\
  db? : DBNAME
  \where
  sessionPhase = spProcessing \\
  toolPhase = tpNone \\
  %
  activeDb' = \{ db? \} \\
  %
  sessionPhase' = sessionPhase \\
  sessionId' = sessionId \\
  permMode' = permMode \\
  toolPhase' = toolPhase \\
  currentTool' = currentTool \\
  toolsThisTurn' = toolsThisTurn \\
  activeSubagents' = activeSubagents \\
  context' = context \\
  pendingPrompt' = pendingPrompt \\
  pendingResponse' = pendingResponse \\
  codebaseIndexed' = codebaseIndexed \\
  capturedUrls' = capturedUrls \\
  compactionSaved' = compactionSaved \\
  cfgSessionSync' = cfgSessionSync \\
  cfgWebFetch' = cfgWebFetch \\
  cfgCompaction' = cfgCompaction
\end{schema}

\subsection{SearchKnowledge}

Claude calls \texttt{quarry find}. This is a read-only query that
injects search results into context. Modeled to show that search
results flow into the conversation.

\begin{schema}{SearchKnowledge}
  \Delta State \\
  results? : CONTEXTCHUNK
  \where
  sessionPhase = spProcessing \\
  toolPhase = tpNone \\
  %
  context' = (1 \upto maxContext) \dres (context \cat \langle results? \rangle) \\
  %
  sessionPhase' = sessionPhase \\
  sessionId' = sessionId \\
  permMode' = permMode \\
  toolPhase' = toolPhase \\
  currentTool' = currentTool \\
  toolsThisTurn' = toolsThisTurn \\
  activeSubagents' = activeSubagents \\
  pendingPrompt' = pendingPrompt \\
  pendingResponse' = pendingResponse \\
  codebaseIndexed' = codebaseIndexed \\
  activeDb' = activeDb \\
  capturedUrls' = capturedUrls \\
  compactionSaved' = compactionSaved \\
  cfgSessionSync' = cfgSessionSync \\
  cfgWebFetch' = cfgWebFetch \\
  cfgCompaction' = cfgCompaction
\end{schema}

\subsection{InjectKnowledgeHint}

Quarry proactively injects relevant knowledge into context at key
moments: \texttt{SessionStart} (project standards, workflow tips),
\texttt{PostToolUse} on \texttt{Edit|Write} (e.g., ``remember our
changelog standard'', ``update DESIGN.md with an ADR''),
\texttt{PostToolUse} on tool-specific patterns (e.g., ``top hints
for working with Z specs'' when editing \texttt{.tex} files).

These are contextual hints retrieved from quarry collections
(standards, design docs, project conventions) and injected as
\texttt{additionalContext}. The injection is non-blocking and
additive --- it appends to context without altering other state.

\begin{schema}{InjectKnowledgeHint}
  \Delta State \\
  hint? : CONTEXTCHUNK
  \where
  sessionPhase \in \{ spIdle, spProcessing \} \\
  %
  context' = (1 \upto maxContext) \dres (context \cat \langle hint? \rangle) \\
  %
  sessionPhase' = sessionPhase \\
  sessionId' = sessionId \\
  permMode' = permMode \\
  toolPhase' = toolPhase \\
  currentTool' = currentTool \\
  toolsThisTurn' = toolsThisTurn \\
  activeSubagents' = activeSubagents \\
  pendingPrompt' = pendingPrompt \\
  pendingResponse' = pendingResponse \\
  codebaseIndexed' = codebaseIndexed \\
  activeDb' = activeDb \\
  capturedUrls' = capturedUrls \\
  compactionSaved' = compactionSaved \\
  cfgSessionSync' = cfgSessionSync \\
  cfgWebFetch' = cfgWebFetch \\
  cfgCompaction' = cfgCompaction
\end{schema}

This operation fires in both \texttt{spIdle} (SessionStart context
injection) and \texttt{spProcessing} (PostToolUse context injection).
The hint content is determined by quarry's semantic search over
registered collections --- the Z spec abstracts this as an opaque
\texttt{CONTEXTCHUNK} input.

Example injection points:

\begin{itemize}
  \item \texttt{SessionStart}: ``This project follows punt-kit
    standards. Key conventions: ruff, mypy, pytest, CHANGELOG.md
    under [Unreleased].''
  \item \texttt{PostToolUse} on \texttt{Edit} of \texttt{CHANGELOG.md}:
    ``Keep a Changelog format. Entries under [Unreleased] with
    Added/Changed/Fixed sections.''
  \item \texttt{PostToolUse} on \texttt{Edit} of \texttt{*.tex}:
    ``Z spec conventions: use $\backslash$quad\~{} for indentation,
    ZBOOL for booleans, bounded sets for ProB animation.''
  \item \texttt{PostToolUse} on \texttt{Bash} with git commands:
    ``Session close protocol: git status, git add, bd sync,
    git commit, git push.''
\end{itemize}

% ============================================================
\section{Subagent Operations}
% ============================================================

\begin{schema}{SpawnSubagent}
  QuarryFrame \\
  agentId? : AGENTID
  \where
  sessionPhase = spProcessing \\
  agentId? \notin activeSubagents \\
  \# activeSubagents < maxSubagents \\
  activeSubagents' = activeSubagents \cup \{ agentId? \} \\
  sessionPhase' = sessionPhase \\
  sessionId' = sessionId \\
  permMode' = permMode \\
  toolPhase' = toolPhase \\
  currentTool' = currentTool \\
  toolsThisTurn' = toolsThisTurn \\
  context' = context \\
  pendingPrompt' = pendingPrompt \\
  pendingResponse' = pendingResponse
\end{schema}

\begin{schema}{StopSubagent}
  QuarryFrame \\
  agentId? : AGENTID \\
  hookResult? : HookDecision \\
  resultContext? : CONTEXTCHUNK
  \where
  sessionPhase = spProcessing \\
  agentId? \in activeSubagents \\
  hookResult? = hdAllow \\
  activeSubagents' = activeSubagents \setminus \{ agentId? \} \\
  context' = (1 \upto maxContext) \dres (context \cat \langle resultContext? \rangle) \\
  sessionPhase' = sessionPhase \\
  sessionId' = sessionId \\
  permMode' = permMode \\
  toolPhase' = toolPhase \\
  currentTool' = currentTool \\
  toolsThisTurn' = toolsThisTurn \\
  pendingPrompt' = pendingPrompt \\
  pendingResponse' = pendingResponse
\end{schema}

% ============================================================
\section{Response and Stop}
% ============================================================

\begin{schema}{BeginResponse}
  QuarryFrame \\
  response? : RESPONSE
  \where
  sessionPhase = spProcessing \\
  toolPhase = tpNone \\
  activeSubagents = \emptyset \\
  sessionPhase' = spResponding \\
  pendingResponse' = \{ response? \} \\
  pendingPrompt' = \emptyset \\
  sessionId' = sessionId \\
  permMode' = permMode \\
  toolPhase' = tpNone \\
  currentTool' = \emptyset \\
  toolsThisTurn' = toolsThisTurn \\
  activeSubagents' = activeSubagents \\
  context' = context
\end{schema}

\begin{schema}{FinishResponse}
  QuarryFrame \\
  hookResult? : HookDecision \\
  responseContext? : CONTEXTCHUNK
  \where
  sessionPhase = spResponding \\
  hookResult? = hdAllow \\
  sessionPhase' = spIdle \\
  pendingResponse' = \emptyset \\
  context' = (1 \upto maxContext) \dres (context \cat \langle responseContext? \rangle) \\
  toolsThisTurn' = 0 \\
  sessionId' = sessionId \\
  permMode' = permMode \\
  toolPhase' = tpNone \\
  currentTool' = \emptyset \\
  activeSubagents' = activeSubagents \\
  pendingPrompt' = \emptyset
\end{schema}

\begin{schema}{ForceContinue}
  QuarryFrame \\
  hookResult? : HookDecision
  \where
  sessionPhase = spResponding \\
  hookResult? = hdBlock \\
  sessionPhase' = spProcessing \\
  pendingResponse' = \emptyset \\
  sessionId' = sessionId \\
  permMode' = permMode \\
  toolPhase' = tpNone \\
  currentTool' = \emptyset \\
  toolsThisTurn' = toolsThisTurn \\
  activeSubagents' = activeSubagents \\
  context' = context \\
  pendingPrompt' = pendingPrompt
\end{schema}

% ============================================================
\section{Context Management --- Compaction Safety}
% ============================================================

This is the critical section for quarry. Context compaction destroys
conversation history. The \texttt{PreCompact} hook must capture
session notes \emph{before} the context is shrunk.

\subsection{CaptureBeforeCompact}

\texttt{PreCompact} hook fires. If compaction capture is enabled,
quarry reads the transcript and ingests it as session notes into
the \texttt{session-notes} collection. This must happen before
\texttt{CompactContext}.

\begin{schema}{CaptureBeforeCompact}
  \Delta State
  \where
  sessionPhase = spIdle \\
  cfgCompaction = ztrue \\
  \# context > 0 \\
  %
  compactionSaved' = ztrue \\
  %
  sessionPhase' = sessionPhase \\
  sessionId' = sessionId \\
  permMode' = permMode \\
  toolPhase' = toolPhase \\
  currentTool' = currentTool \\
  toolsThisTurn' = toolsThisTurn \\
  activeSubagents' = activeSubagents \\
  context' = context \\
  pendingPrompt' = pendingPrompt \\
  pendingResponse' = pendingResponse \\
  codebaseIndexed' = codebaseIndexed \\
  activeDb' = activeDb \\
  capturedUrls' = capturedUrls \\
  cfgSessionSync' = cfgSessionSync \\
  cfgWebFetch' = cfgWebFetch \\
  cfgCompaction' = cfgCompaction
\end{schema}

\subsection{SkipCaptureBeforeCompact}

Compaction capture is disabled. The \texttt{PreCompact} hook is a
no-op. Session notes are not saved.

\begin{schema}{SkipCaptureBeforeCompact}
  \Xi State
  \where
  sessionPhase = spIdle \\
  cfgCompaction = zfalse
\end{schema}

\subsection{CompactContext}

Context is shrunk. The \texttt{compactionSaved} flag records whether
notes were captured. After compaction it resets to \texttt{zfalse}
for the next potential compaction cycle.

\begin{schema}{CompactContext}
  \Delta State \\
  newLen? : \nat
  \where
  sessionPhase = spIdle \\
  \# context > 0 \\
  newLen? < \# context \\
  newLen? \leq maxContext \\
  %
  context' = (1 \upto newLen?) \dres context \\
  compactionSaved' = zfalse \\
  %
  sessionPhase' = sessionPhase \\
  sessionId' = sessionId \\
  permMode' = permMode \\
  toolPhase' = toolPhase \\
  currentTool' = currentTool \\
  toolsThisTurn' = toolsThisTurn \\
  activeSubagents' = activeSubagents \\
  pendingPrompt' = pendingPrompt \\
  pendingResponse' = pendingResponse \\
  codebaseIndexed' = codebaseIndexed \\
  activeDb' = activeDb \\
  capturedUrls' = capturedUrls \\
  cfgSessionSync' = cfgSessionSync \\
  cfgWebFetch' = cfgWebFetch \\
  cfgCompaction' = cfgCompaction
\end{schema}

% ============================================================
\section{Session End}
% ============================================================

Every resource the session held is released; the session identifier is
kept, because it is the record of which session ended. See
\texttt{claude-code.tex} for the discussion.

\begin{schema}{EndSession}
  \Delta State \\
  reason? : SessionEndReason
  \where
  sessionPhase = spIdle \\
  activeSubagents = \emptyset \\
  sessionPhase' = spEnded \\
  sessionId' = sessionId \\
  toolPhase' = tpNone \\
  currentTool' = \emptyset \\
  toolsThisTurn' = 0 \\
  activeSubagents' = \emptyset \\
  context' = \langle \rangle \\
  pendingPrompt' = \emptyset \\
  pendingResponse' = \emptyset \\
  permMode' = permMode \\
  codebaseIndexed' = zfalse \\
  activeDb' = \emptyset \\
  capturedUrls' = \emptyset \\
  compactionSaved' = zfalse \\
  cfgSessionSync' = cfgSessionSync \\
  cfgWebFetch' = cfgWebFetch \\
  cfgCompaction' = cfgCompaction
\end{schema}

\subsection{SessionOver}

The declared halt: enabled exactly in \texttt{spEnded}, and it changes
nothing. Stating the intended termination as an operation is what makes
the absence of others meaningful --- a model checker cannot otherwise
distinguish a machine that has finished from one that is stuck, so with
this schema present every remaining deadlock is unintended by
construction.

\begin{schema}{SessionOver}
  \Xi State
  \where
  sessionPhase = spEnded
\end{schema}


% ============================================================
\section{System Invariants}
% ============================================================

\begin{enumerate}
  \item \textbf{Codebase indexed when sync enabled}: if
    \texttt{cfgSessionSync = ztrue} and the session started via
    \texttt{StartSessionWithSync}, then
    \texttt{codebaseIndexed = ztrue} for the duration of the session.

  \item \textbf{Capture config is immutable during session}: the
    \texttt{cfg*} flags are established by the \texttt{SessionStart}
    operation that begins the session, and preserved unchanged by every
    operation thereafter. They cannot change mid-session.

  \item \textbf{WebFetch deduplication}: \texttt{CaptureWebFetch}
    requires \texttt{url? $\notin$ capturedUrls}. No URL is ingested
    twice in the same session.

  \item \textbf{Active database consistency}: at most one database
    is active. \texttt{SwitchDatabase} replaces, does not accumulate.

  \item \textbf{Compaction flag lifecycle}:
    \texttt{CaptureBeforeCompact} sets \texttt{compactionSaved = ztrue}.
    \texttt{CompactContext} resets it to \texttt{zfalse}. This tracks
    whether the most recent compaction was preceded by a capture.

  \item \textbf{Clean session boundaries}: all quarry state
    (indexed, activeDb, captures, compaction flag) is cleared on
    \texttt{EndSession}.
\end{enumerate}

\textbf{Note on the compaction safety property:} The Z spec does
\emph{not} enforce that \texttt{CaptureBeforeCompact} must precede
\texttt{CompactContext} --- it is possible to compact without
capturing (the operations are independent). This reflects reality:
Claude Code fires \texttt{PreCompact} before compaction, but the
hook is non-blocking. The \texttt{compactionSaved} flag lets us
\emph{observe} whether the capture occurred, not \emph{enforce} it.
Enforcement is the hook system's responsibility.

% ============================================================
\section{Hook Event Map}
% ============================================================

\begin{tabular}{lllll}
\hline
\textbf{Hook} & \textbf{Matcher} & \textbf{Operation} & \textbf{Quarry Effect} & \textbf{Blocks?} \\
\hline
SessionStart & * & StartSessionWithSync/NoSync & Auto-register + sync & No \\
SessionStart & * & InjectKnowledgeHint & Project standards & No \\
PostToolUse & WebFetch & CaptureWebFetch & Ingest URL & No \\
PostToolUse & Edit|Write & InjectKnowledgeHint & File-type hints & No \\
PostToolUse & Bash & InjectKnowledgeHint & Workflow hints & No \\
PostToolUse & quarry\_\_use & SwitchDatabase & Change active DB & No \\
PostToolUse & quarry\_\_find & SearchKnowledge & Inject results & No \\
PostToolUse & quarry\_\_* & (suppress output) & Format display & No \\
PreCompact & * & CaptureBeforeCompact & Save session notes & No \\
SessionEnd & * & EndSession & Clear state & No \\
\hline
\end{tabular}

% ============================================================
\section{Implementation Validation}
% ============================================================

This section cross-references the Z specification against the
current quarry implementation (\texttt{hooks/hooks.json},
\texttt{hooks/*.sh}, \texttt{src/quarry/hooks.py}).

\subsection{Critical Finding: Python Handlers Never Invoked}

The quarry hook system has a \textbf{wiring gap}. Three Python
handlers were implemented in \texttt{src/quarry/hooks.py} but are
never called by the shell scripts registered in \texttt{hooks.json}:

\begin{itemize}
  \item \texttt{handle\_session\_start()} --- auto-registers directory
    and runs sync. \textbf{Never invoked.} The shell script
    \texttt{session-start.sh} only deploys commands and sets
    permissions --- it does not call \texttt{quarry hooks session-start}.
  \item \texttt{handle\_post\_web\_fetch()} --- ingests fetched URLs
    into \texttt{web-captures} collection. \textbf{Not registered} in
    \texttt{hooks.json} at all. No \texttt{PostToolUse} hook on
    \texttt{WebFetch} exists.
  \item \texttt{handle\_pre\_compact()} --- captures transcript as
    session notes. \textbf{Not registered} in \texttt{hooks.json}.
    No \texttt{PreCompact} hook exists.
\end{itemize}

\textbf{Root cause:} The implementation was split between shell
scripts (setup/formatting) and Python handlers (knowledge capture).
The Python handlers were written but never wired into the hook
dispatch chain. The shell scripts do not call back to Python.

\subsection{Built and Correct}

\begin{tabular}{llll}
\hline
\textbf{Component} & \textbf{Hook} & \textbf{File} & \textbf{Status} \\
\hline
Command deployment & SessionStart & session-start.sh & Correct \\
MCP tool permissions & SessionStart & session-start.sh & Correct \\
Output formatting & PostToolUse (quarry) & suppress-output.sh & Correct \\
Hook config parsing & (library) & hooks.py & Correct \\
\hline
\end{tabular}

\subsection{Implemented but Not Wired}

\begin{tabular}{llll}
\hline
\textbf{Z Operation} & \textbf{Handler} & \textbf{hooks.json} & \textbf{Fix} \\
\hline
StartSessionWithSync & handle\_session\_start() & Registered but & Wire shell to Python \\
 & & shell doesn't call it & \\
CaptureWebFetch & handle\_post\_web\_fetch() & Not registered & Add PostToolUse WebFetch \\
CaptureBeforeCompact & handle\_pre\_compact() & Not registered & Add PreCompact entry \\
\hline
\end{tabular}

\textbf{Fix pattern:} Each shell script needs to call
\texttt{quarry hooks <subcommand> < /dev/stdin} to delegate to the
Python handler. Alternatively, register the Python CLI directly as
the hook command.

\subsection{Not Yet Built}

\begin{tabular}{llll}
\hline
\textbf{Z Operation} & \textbf{Hook Event} & \textbf{Effect} & \textbf{Effort} \\
\hline
InjectKnowledgeHint & SessionStart & Project standards & Medium \\
InjectKnowledgeHint & PostToolUse Edit|Write & File-type hints & Medium \\
InjectKnowledgeHint & PostToolUse Bash & Workflow hints & Medium \\
EndSession cleanup & SessionEnd & Clear session state & Low \\
\hline
\end{tabular}

\textbf{Knowledge hint injection} is the most significant new
capability. It requires:

\begin{enumerate}
  \item A hint resolution engine that maps the current activity
    (tool name, file path, command) to relevant quarry collections.
  \item A PostToolUse hook on \texttt{Edit|Write} that detects the
    file being edited and searches quarry for relevant standards.
  \item A SessionStart hook that injects project-level conventions.
  \item Careful scoping: hints must be short (1--3 sentences) to
    avoid flooding context. Quarry's semantic search ranks by
    relevance --- the top result should be the hint.
\end{enumerate}

\subsection{Coverage Summary}

\begin{tabular}{lll}
\hline
\textbf{Z Operation} & \textbf{Implementation} & \textbf{Status} \\
\hline
StartSessionWithSync & hooks.py (handler exists) & \textbf{Not wired} \\
StartSessionNoSync & hooks.py (config check) & \textbf{Not wired} \\
SubmitPrompt / RejectPrompt & (base) & Correct \\
Tool call operations & (base) & Correct \\
CaptureWebFetch & hooks.py (handler exists) & \textbf{Not registered} \\
SwitchDatabase & MCP tool & Correct \\
SearchKnowledge & MCP tool & Correct \\
InjectKnowledgeHint & --- & \textbf{Not built} \\
CaptureBeforeCompact & hooks.py (handler exists) & \textbf{Not registered} \\
CompactContext & (base) & Correct \\
EndSession & --- & \textbf{Not built} \\
\hline
\end{tabular}

\end{document}
